\section{Background: FTS-Diffusion}
\label{sec:background}

% Target: 1 page (~800 words). Recap enough that a reader can follow §4–§6
% without reading the original paper, but do not redo the original derivations.

\subsection{Architecture}
\label{sec:background:arch}

% Three modules (per Huang et al., Fig.~3):
% 1. Pattern Recognition Module (PRM) — Scale-Invariant Subsequence Clustering
%    (SISC): variable-length segmentation + DTW-based k-means clustering. K
%    pattern centroids.
% 2. Pattern Generation Module (PGM) — pattern-conditioned diffusion (DDPM)
%    on a fixed-length latent, paired with a Scaling AutoEncoder (SAE) that
%    handles the variable-length ↔ fixed-length mapping. Trained jointly to
%    minimize reconstruction MSE + denoising ELBO.
% 3. Pattern Evolution Module (PEM) — a feed-forward network mapping a state
%    (p_t, α_t, β_t) to predicted (p_{t+1}, α_{t+1}, β_{t+1}). Trained with
%    cross-entropy on pattern + length, MSE/L1 on magnitude.

\subsection{Sampling pipeline}
\label{sec:background:sampling}

% Autoregressive: start from an initial segment from real data, then loop:
% (a) PEM predicts next state, (b) PGM generates segment from state +
% pattern centroid, (c) anchor: segment ← segment - segment[0] + ts[-1],
% (d) append.
%
% In the released code, the PEM next-state prediction uses argmax for the
% pattern and length heads (deterministic) — we revisit this in §4.

\subsection{Evaluation protocol}
\label{sec:background:eval}

% The paper evaluates synthetic series in three ways:
% 1. Stylized facts: visual + qualitative — heavy tails in returns, decaying
%    autocorrelation of |returns|.
% 2. Distribution tests: KS and AD on the marginal return distribution; higher
%    is better (statistics are capped, see Table 1 in original paper).
% 3. Downstream prediction: an LSTM predicts the next day's price from the
%    previous 64 days. Two settings:
%    - TMTR (Train on Mixture, Test on Real): vary the real/synthetic mix of
%      a fixed-size training set; measure MAPE on a held-out real test set.
%    - TATR (Train on Augmentation, Test on Real): grow a training set
%      starting from a small real seed by appending synthetic data; measure
%      MAPE as the augmented size grows.
% The paper reports KS/AD ~0.32/0.12 on three assets and MAPE reductions up to
% 17.9% on TATR with 100 years of synthetic data appended.

\todo{Decide whether to include a redrawn architecture figure (Fig.~3 from
original paper) or just describe in text. If figure: place after
\S\ref{sec:background:arch}; budget another quarter page.}
